\documentclass[11pt]{article}
% ======================================================================
%  weekpreamble.tex — shared preamble for the weekly course summaries (EN)
%  Concise, readable, intuition-first. Same box style as the main doc.
%  Compiles with pdflatex.
% ======================================================================
\usepackage[utf8]{inputenc}
\usepackage[T1]{fontenc}
\usepackage[english]{babel}
\usepackage[a4paper,margin=2.2cm]{geometry}
\usepackage{amsmath,amssymb,amsthm,mathtools}
\usepackage{bm}
\usepackage{enumitem}
\usepackage{xcolor}
\usepackage[most]{tcolorbox}
\usepackage{microtype}

\emergencystretch=3em
\allowdisplaybreaks[1]
\setlength{\parindent}{0pt}
\setlength{\parskip}{0.45em}

% ---------- math macros (same names as the main document) ----------
\newcommand{\E}{\mathbb{E}}
\DeclareMathOperator{\Var}{Var}
\DeclareMathOperator{\Cov}{Cov}
\DeclareMathOperator{\Corr}{Corr}
\DeclareMathOperator*{\argmin}{arg\,min}
\DeclareMathOperator{\MSE}{MSE}
\DeclareMathOperator{\trace}{tr}
\newcommand{\Reals}{\mathbb{R}}
\newcommand{\Ints}{\mathbb{Z}}
\newcommand{\Comp}{\mathbb{C}}
\newcommand{\Nat}{\mathbb{N}}
\newcommand{\Prob}{\mathbb{P}}
\newcommand{\Tr}{^{\mathsf{T}}}
\newcommand{\Herm}{^{\mathsf{H}}}
\newcommand{\conj}[1]{\overline{#1}}
\newcommand{\abs}[1]{\left\lvert #1 \right\rvert}
\newcommand{\norm}[1]{\left\lVert #1 \right\rVert}
\newcommand{\ind}{\mathbf{1}}
\newcommand{\eps}{\varepsilon}
\newcommand{\diff}{\nabla}
\newcommand{\acvs}{\gamma}
\newcommand{\acf}{\rho}
\newcommand{\pacf}{\alpha}
\newcommand{\sdf}{\mathcal{S}}
\newcommand{\filt}{\mathcal{L}}
\newcommand{\Bsh}{B}
\newcommand{\ms}{\;\overset{\mathrm{ms}}{=}\;}
\newcommand{\toprob}{\xrightarrow{P}}
\newcommand{\todist}{\xrightarrow{d}}
\newcommand{\iid}{\overset{\text{iid}}{\sim}}

\setlist[itemize]{leftmargin=1.5em,topsep=2pt,itemsep=2pt}
\setlist[enumerate]{leftmargin=1.7em,topsep=2pt,itemsep=2pt}

% ---------- compact, titled (unnumbered) boxes ----------
\tcbset{wkbase/.style={enhanced,breakable,fonttitle=\bfseries,coltitle=white,
  boxrule=0.6pt,arc=1.2mm,left=2.2mm,right=2.2mm,top=1.1mm,bottom=1.1mm,
  before skip=6pt,after skip=6pt,
  attach boxed title to top left={yshift=-3.2mm,xshift=2.4mm},
  boxed title style={boxrule=0pt,arc=0.5mm}}}

\newtcolorbox{defn}[1]{wkbase, colback=blue!4!white, colframe=blue!58!black,
  colbacktitle=blue!58!black, title={Definition --- #1}, top=3mm}
\newtcolorbox{result}[1]{wkbase, colback=red!4!white, colframe=red!60!black,
  colbacktitle=red!60!black, title={#1}, top=3mm}
\newtcolorbox{intu}[1][Intuition]{wkbase, colback=yellow!12!white,
  colframe=orange!72!black, colbacktitle=orange!72!black, coltitle=white,
  title={#1}, top=3mm}
\newtcolorbox{retenir}[1][Key takeaways --- the week's reflexes]{wkbase,
  colback=green!5!white, colframe=green!48!black, colbacktitle=green!48!black,
  title={#1}, top=3mm}
\newtcolorbox{exmpl}[1][Worked example]{wkbase, colback=black!3!white,
  colframe=black!55!white, colbacktitle=black!55!white, title={#1}, top=3mm}

% ---------- week header ----------
\newcommand{\weekheader}[4]{%
  {\small\sffamily\bfseries\color{black!60} TIME SERIES~~$\cdot$~~MATH-342, EPFL}\par
  \vspace{2pt}
  {\LARGE\bfseries Week #1 --- #3\par}
  \vspace{1pt}
  {\itshape\color{black!55} #2}\par
  \vspace{6pt}
  \begin{tcolorbox}[enhanced,colback=blue!3!white,colframe=blue!45!black,
    arc=1.4mm,boxrule=0.7pt,left=3mm,right=3mm,top=2mm,bottom=2mm,before skip=2pt,after skip=10pt]
    \textbf{The big picture.}~ #4
  \end{tcolorbox}
}

\begin{document}
\weekheader{09}{23 April 2026}{Multivariate Time Series \& VAR}{Move to vectors: matrix autocovariance (the lag is in the FIRST argument!), spectral matrix, coherence, and the VAR model with its companion form.}

We now observe several series at once: a basket of stocks (Apple, Nasdaq-100, S\&P 500), or the position and temperature of an ocean drifter. Each $X_t$ becomes a \emph{vector} in $\Reals^d$. Two questions: how do we \emph{describe} the second-order structure (now matrix-valued, with genuine \textbf{cross}-terms between components), and how do we \emph{model} it (answer: the VAR). The recurring trick: \emph{every VAR($p$) is a VAR(1) in disguise}, so the whole theory reduces to a single matrix.

% ======================================================================
\section*{Key definitions}

\begin{defn}{Multivariate series \& second-order stationarity}
$\{X_t\}$ with $X_t=(X_t^{(1)},\dots,X_t^{(d)})\Tr\in\Reals^d$, each marginal a univariate series. It is \textbf{(weakly) second-order stationary} if the mean is constant, $\trace(\Var(X_t))<\infty$, and
\[
\Cov(X_{t+\tau},X_t)\ \text{depends only on }\tau .
\]
Equivalently: each marginal is stationary \emph{and} every pair has a shift-invariant cross-covariance.
\end{defn}

\begin{defn}{Matrix autocovariance sequence (ACVS) \& cross-correlation}
\[
\Gamma_\tau=\Cov(X_{t+\tau},X_t),\qquad
\acvs_\tau^{(j,k)}=\Cov\!\big(X_{t+\tau}^{(j)},X_t^{(k)}\big).
\]
\textbf{The lag $\tau$ sits in the FIRST argument} (irrelevant in the scalar case, crucial here). Diagonal ($j=k$): the usual marginal ACVS. Off-diagonal ($j\ne k$): the \textbf{cross-covariance}, which encodes lead--lag relations between channels. The \textbf{cross-correlation sequence} (CCS) is
\[
\acf_\tau^{(j,k)}=\frac{\acvs_\tau^{(j,k)}}{\sqrt{\acvs_0^{(j,j)}\,\acvs_0^{(k,k)}}},\qquad \acf_\tau^{(j,k)}=\acf_{-\tau}^{(k,j)} .
\]
\end{defn}

\begin{defn}{Multivariate white noise $\mathrm{WN}(0,\Sigma)$}
$\{\eps_t\}\sim\mathrm{WN}(0,\Sigma)$ iff it is stationary, mean zero, with $\Gamma_\tau^{(\eps)}=\Sigma$ at $\tau=0$ and $0$ otherwise. Components are uncorrelated \textbf{across time} but may be correlated \textbf{contemporaneously} (off-diagonal entries of $\Sigma$). One sometimes assumes $\Sigma$ diagonal --- an \emph{extra} assumption, not part of the definition.
\end{defn}

\begin{defn}{Spectral density matrix}
If all cross-covariances are summable ($\ell^1$),
\[
S(f)=\sum_{\tau\in\Ints}\Gamma_\tau\,e^{-2\pi i\tau f},\qquad
\Gamma_\tau=\int_{-1/2}^{1/2}S(f)\,e^{2\pi if\tau}\,df,\qquad f\in[-\tfrac12,\tfrac12].
\]
Diagonal: the usual (real) spectral density $\sdf_j(f)$. Off-diagonal: the \textbf{cross-spectrum} $S_{j,k}(f)\in\Comp$ (complex in general).
\end{defn}

\begin{defn}{Coherence}
\[
r_{j,k}(f)=\frac{S_{j,k}(f)}{\sqrt{S_{j,j}(f)\,S_{k,k}(f)}}\in\Comp .
\]
The frequency-by-frequency ``correlation'' between channels (literally the correlation of the increments $dZ_j,dZ_k$ below). Report its \textbf{magnitude} $\abs{r_{j,k}(f)}$ (strength of association, $\abs{r}^2\in[0,1]$) and its \textbf{phase} $\arg r_{j,k}(f)$ (the ``sign'' / lead--lag angle).
\end{defn}

\begin{defn}{VAR($p$) \& companion form}
With $\{\eps_t\}\sim\mathrm{WN}(0,\Sigma)$ and $\Phi_p\ne0$,
\[
X_t=\Phi_1 X_{t-1}+\cdots+\Phi_p X_{t-p}+\eps_t,\qquad \Phi(\Bsh)X_t=\eps_t,\quad \Phi(z)=I_d-\sum_{j=1}^p\Phi_j z^j .
\]
Stacking $Y_t=(X_t\Tr,\dots,X_{t-p+1}\Tr)\Tr\in\Reals^{dp}$ turns it into a VAR(1), $Y_t=F\,Y_{t-1}+U_t$, with
\[
F=\begin{pmatrix}
\Phi_1 & \Phi_2 & \cdots & \Phi_{p-1} & \Phi_p\\
I_d & 0 & \cdots & 0 & 0\\
0 & I_d & & & \vdots\\
\vdots & & \ddots & & \vdots\\
0 & \cdots & & I_d & 0
\end{pmatrix},
\quad U_t=\begin{pmatrix}\eps_t\\0\\\vdots\\0\end{pmatrix}.
\]
Recover the original process via $X_t=G\,Y_t$ with $G=(I_d\ 0\ \cdots\ 0)$. The off-diagonal entries of the $\Phi_j$ \emph{are} the cross-component dynamics.
\end{defn}

% ======================================================================
\section*{Key results \& formulas}

\begin{result}{Symmetry of the ACVS and the spectrum}
\[
\acvs_\tau^{(j,k)}=\acvs_{-\tau}^{(k,j)}\ \Longleftrightarrow\ \Gamma_{-\tau}=\Gamma_\tau\Tr,
\qquad S(f)=S(f)\Herm=\conj{S(-f)\Tr},\quad S(f)\succeq0 .
\]
$\{\Gamma_\tau\}$ is positive semi-definite: $\sum_{j,k}a_j\Tr\Gamma_{k-j}a_k\ge0$ (it is a variance). \textbf{Warning:} off-diagonal, the ACVS is \emph{not} even in $\tau$; you must swap the indices $(j,k)\!\to\!(k,j)$ when you flip the sign of $\tau$. This asymmetry is exactly what encodes lead--lag.
\end{result}

\begin{result}{Theorem --- Multivariate spectral representation}
There is a vector orthogonal-increment process $\{Z(f)\}$ with
\[
\begin{aligned}
X_t&=\mu+\int_{-1/2}^{1/2}e^{2\pi ift}\,dZ(f),\\
\E[dZ(f)]&=0,\quad \Var\big(dZ(f)\big)=S(f)\,df,\quad \Cov\big(dZ(f),dZ(f')\big)=0\ (f\ne f').
\end{aligned}
\]
\emph{Idea:} the vector analogue of Cram\'er. So a stationary vector process is a sum of uncorrelated random sinusoids whose (matrix) variance is $S(f)\,df$; coherence is literally the correlation of $dZ_j(f),dZ_k(f)$.
\end{result}

\begin{result}{Stability of a VAR($p$)}
Stable iff the roots of the reverse characteristic polynomial $\det(I_d-\Phi_1 z-\cdots-\Phi_p z^p)=0$ \emph{all} lie outside the unit circle ($\abs{z}>1$); equivalently iff $F$ has all eigenvalues $\abs{\lambda}<1$. Stability $\Rightarrow$ stationarity (not conversely in general). \textbf{VAR(1) shortcut:} just check $\abs{\lambda_i(\Phi_1)}<1$ (spectral radius $<1$); for triangular $\Phi_1$ these are the diagonal entries. Sanity check: $\det(\cdot)\big|_{z=0}=\det I_d=1$.
\end{result}

\begin{result}{MA($\infty$) and covariance of a stable VAR(1)}
\[
X_t=\sum_{j\ge0}\Phi_1^{\,j}\,\eps_{t-j},\qquad
\Gamma_\tau=\sum_{k\ge0}\Phi_1^{\,k+\tau}\,\Sigma\,(\Phi_1^{\,k})\Tr\quad(\tau\ge0).
\]
At $\tau=0$ this is the discrete Lyapunov equation $\Gamma_0=\Phi_1\Gamma_0\Phi_1\Tr+\Sigma$. For a VAR($p$): same formula with $F,\Sigma_U$ (top-left block $\Sigma$), then project $\Gamma_\tau^{(X)}=G\big(\sum_k F^{k+\tau}\Sigma_U(F^k)\Tr\big)G\Tr$.
\end{result}

\begin{result}{Theorem --- Yule--Walker equations (VAR)}
For $\tau\ge0$,
\[
\Gamma_\tau=\Phi_1\Gamma_{\tau-1}+\cdots+\Phi_p\Gamma_{\tau-p}+\delta_{\tau,0}\,\Sigma .
\]
\emph{Idea:} multiply the VAR equation on the right by $X_t\Tr$ and take expectations; the noise term contributes $\Sigma$ only at $\tau=0$. Solve the first $p+1$ equations for $\Gamma_0,\dots,\Gamma_p$, then recurse for all $\tau>p$. Run in reverse, it \emph{estimates} the $\Phi_j$ from sample covariances (VAR fitting).
\end{result}

\begin{result}{Two signatures: common noise \& pure delay}
\emph{Common noise.} If $Y_t=X_t+\phi_t\ind$ (the \emph{same} scalar noise in both channels), then $\Gamma_\tau^{(Y)}=\Gamma_\tau^{(X)}+\sigma_\phi^2\,\delta_{\tau,0}\,\ind\ind\Tr$: correlation added \emph{only} at lag $0$.\\
\emph{Pure delay.} If $X_t=(Y_t,Y_{t-\nu})\Tr$, then $S_{1,2}(f)=\sdf_Y(f)\,e^{2\pi if\nu}$, hence coherence of magnitude $1$ and phase $2\pi\nu f$ (\emph{linear} in $f$).
\end{result}

% ======================================================================
\section*{Intuition}

\begin{intu}[Why the lag lives in the first argument]
Scalar: $\acvs_\tau=\acvs_{-\tau}$, so argument order is irrelevant. Multivariate: it becomes \emph{crucial} --- it fixes the convention $\Gamma_{-\tau}=\Gamma_\tau\Tr$. Diagonals stay even, but a cross-covariance $\acvs_\tau^{(1,2)}$ (channel $1$ leading channel $2$ by $\tau$) has no reason to equal $\acvs_{-\tau}^{(1,2)}$. Classic error: forgetting to swap $(j,k)$ when you flip $\tau$.
\end{intu}

\begin{intu}[The noise only shows up at $\tau=0$]
$X_t$ is built from $\eps_t,\eps_{t-1},\dots$ (the causal MA($\infty$)). For $\tau>0$, $\eps_{t+\tau}$ lies in the \emph{future} of $X_t$, hence uncorrelated: $\Cov(\eps_{t+\tau},X_t)=0$. Only $\tau=0$ gives $\Cov(\eps_t,X_t)=\Sigma$. Exactly the scalar Yule--Walker mechanism, transposed to matrices.
\end{intu}

\begin{intu}[``Everything is a VAR(1)'']
The companion form collapses any VAR($p$) into a single VAR(1) on $F$. So you prove the VAR(1) theory \emph{once} (stability, MA($\infty$), covariances), then reuse it for every $p$ via $X_t=GY_t$. The top block row of $F$ carries the real equation; the lower $I_d$ blocks just shift the stack.
\end{intu}

\begin{intu}[Delay $\leftrightarrow$ linear phase]
A shift of $\nu$ samples in time $=$ a factor $e^{\pm2\pi i\nu f}$ in the spectrum $=$ a \emph{linear} phase $\mp2\pi\nu f$. A coherence of magnitude $1$ whose phase grows proportionally to $f$ (slope $2\pi\nu$) is the textbook signature of a lead/lag of $\nu$ samples between two channels.
\end{intu}

\begin{intu}[Where the channel couplings live]
The \emph{off-diagonal} entries of the $\Phi_j$ \emph{and} of $\Sigma$ are the only links between components. If $\Phi_1$ and $\Sigma$ are diagonal, the VAR(1) decouples into $d$ independent scalar AR(1)'s ($\Gamma_0^{(j,j)}=\sigma_j^2/(1-\phi_j^2)$): cross-spectrum zero, coherence zero everywhere.
\end{intu}

% ======================================================================
\section*{Key takeaways}

\begin{retenir}
\begin{itemize}
  \item \textbf{Lag in the first argument} $\Rightarrow$ $\Gamma_{-\tau}=\Gamma_\tau\Tr$ and $\acvs_\tau^{(j,k)}=\acvs_{-\tau}^{(k,j)}$: always swap indices when flipping $\tau$.
  \item \textbf{Test stability:} roots of $\det(I_d-\sum\Phi_j z^j)$ all $\abs{z}>1$, or eigenvalues of $F$ all $\abs{\lambda}<1$. Sanity check: $\det(\cdot)\big|_{z=0}=1$. For VAR(1), read off $\lambda_i(\Phi_1)$ directly.
  \item \textbf{MA($\infty$) \& covariance:} write $X_t=\sum\Phi_1^j\eps_{t-j}$, get $\Gamma_\tau=\sum_k\Phi_1^{k+\tau}\Sigma(\Phi_1^k)\Tr$, or solve Lyapunov for $\Gamma_0$.
  \item \textbf{Yule--Walker recursion:} $\Gamma_\tau=\sum_i\Phi_i\Gamma_{\tau-i}$ (add $\Sigma$ at $\tau=0$); mind the transpose $\Gamma_{-s}=\Gamma_s\Tr$.
  \item \textbf{Cross-spectrum:} find $\acvs_\tau^{(j,k)}$ (often a \emph{shift} of a known covariance), Fourier-transform, re-index to pull out $e^{\pm2\pi i\nu f}$, read magnitude \& phase, divide by $\sqrt{S_{jj}S_{kk}}$ for coherence.
  \item \textbf{Signature reflexes:} common noise $\to$ correlation at lag $0$ only; pure delay $\to$ coherence $1$, linear phase.
\end{itemize}
\end{retenir}

\end{document}
